\documentclass[11pt]{article}

    \usepackage[T2A]{fontenc}
\usepackage[utf8]{inputenc}
\usepackage[english, russian]{babel}

\usepackage[breakable]{tcolorbox}
    \usepackage{parskip} % Stop auto-indenting (to mimic markdown behaviour)
    
    \usepackage{iftex}
    \ifPDFTeX
    	\usepackage[T1]{fontenc}
    	\usepackage{mathpazo}
    \else
    	\usepackage{fontspec}
    \fi

    % Basic figure setup, for now with no caption control since it's done
    % automatically by Pandoc (which extracts ![](path) syntax from Markdown).
    \usepackage{graphicx}
    % Maintain compatibility with old templates. Remove in nbconvert 6.0
    \let\Oldincludegraphics\includegraphics
    % Ensure that by default, figures have no caption (until we provide a
    % proper Figure object with a Caption API and a way to capture that
    % in the conversion process - todo).
    \usepackage{caption}
    \DeclareCaptionFormat{nocaption}{}
    \captionsetup{format=nocaption,aboveskip=0pt,belowskip=0pt}

    \usepackage{float}
    \floatplacement{figure}{H} % forces figures to be placed at the correct location
    \usepackage{xcolor} % Allow colors to be defined
    \usepackage{enumerate} % Needed for markdown enumerations to work
    \usepackage{geometry} % Used to adjust the document margins
    \usepackage{amsmath} % Equations
    \usepackage{amssymb} % Equations
    \usepackage{textcomp} % defines textquotesingle
    % Hack from http://tex.stackexchange.com/a/47451/13684:
    \AtBeginDocument{%
        \def\PYZsq{\textquotesingle}% Upright quotes in Pygmentized code
    }
    \usepackage{upquote} % Upright quotes for verbatim code
    \usepackage{eurosym} % defines \euro
    \usepackage[mathletters]{ucs} % Extended unicode (utf-8) support
    \usepackage{fancyvrb} % verbatim replacement that allows latex
    \usepackage{grffile} % extends the file name processing of package graphics 
                         % to support a larger range
    \makeatletter % fix for old versions of grffile with XeLaTeX
    \@ifpackagelater{grffile}{2019/11/01}
    {
      % Do nothing on new versions
    }
    {
      \def\Gread@@xetex#1{%
        \IfFileExists{"\Gin@base".bb}%
        {\Gread@eps{\Gin@base.bb}}%
        {\Gread@@xetex@aux#1}%
      }
    }
    \makeatother
    \usepackage[Export]{adjustbox} % Used to constrain images to a maximum size
    \adjustboxset{max size={0.9\linewidth}{0.9\paperheight}}

    % The hyperref package gives us a pdf with properly built
    % internal navigation ('pdf bookmarks' for the table of contents,
    % internal cross-reference links, web links for URLs, etc.)
    \usepackage{hyperref}
    % The default LaTeX title has an obnoxious amount of whitespace. By default,
    % titling removes some of it. It also provides customization options.
    \usepackage{titling}
    \usepackage{longtable} % longtable support required by pandoc >1.10
    \usepackage{booktabs}  % table support for pandoc > 1.12.2
    \usepackage[inline]{enumitem} % IRkernel/repr support (it uses the enumerate* environment)
    \usepackage[normalem]{ulem} % ulem is needed to support strikethroughs (\sout)
                                % normalem makes italics be italics, not underlines
    \usepackage{mathrsfs}
    
\usepackage{wrapfig}
\usepackage[rightcaption]{sidecap}
\providecommand{\keywords}[1]{\textbf{\textit{Keywords:}} #1}

\author{A. Yu. Drozdov}

    
    % Colors for the hyperref package
    \definecolor{urlcolor}{rgb}{0,.145,.698}
    \definecolor{linkcolor}{rgb}{.71,0.21,0.01}
    \definecolor{citecolor}{rgb}{.12,.54,.11}

    % ANSI colors
    \definecolor{ansi-black}{HTML}{3E424D}
    \definecolor{ansi-black-intense}{HTML}{282C36}
    \definecolor{ansi-red}{HTML}{E75C58}
    \definecolor{ansi-red-intense}{HTML}{B22B31}
    \definecolor{ansi-green}{HTML}{00A250}
    \definecolor{ansi-green-intense}{HTML}{007427}
    \definecolor{ansi-yellow}{HTML}{DDB62B}
    \definecolor{ansi-yellow-intense}{HTML}{B27D12}
    \definecolor{ansi-blue}{HTML}{208FFB}
    \definecolor{ansi-blue-intense}{HTML}{0065CA}
    \definecolor{ansi-magenta}{HTML}{D160C4}
    \definecolor{ansi-magenta-intense}{HTML}{A03196}
    \definecolor{ansi-cyan}{HTML}{60C6C8}
    \definecolor{ansi-cyan-intense}{HTML}{258F8F}
    \definecolor{ansi-white}{HTML}{C5C1B4}
    \definecolor{ansi-white-intense}{HTML}{A1A6B2}
    \definecolor{ansi-default-inverse-fg}{HTML}{FFFFFF}
    \definecolor{ansi-default-inverse-bg}{HTML}{000000}

    % common color for the border for error outputs.
    \definecolor{outerrorbackground}{HTML}{FFDFDF}

    % commands and environments needed by pandoc snippets
    % extracted from the output of `pandoc -s`
    \providecommand{\tightlist}{%
      \setlength{\itemsep}{0pt}\setlength{\parskip}{0pt}}
    \DefineVerbatimEnvironment{Highlighting}{Verbatim}{commandchars=\\\{\}}
    % Add ',fontsize=\small' for more characters per line
    \newenvironment{Shaded}{}{}
    \newcommand{\KeywordTok}[1]{\textcolor[rgb]{0.00,0.44,0.13}{\textbf{{#1}}}}
    \newcommand{\DataTypeTok}[1]{\textcolor[rgb]{0.56,0.13,0.00}{{#1}}}
    \newcommand{\DecValTok}[1]{\textcolor[rgb]{0.25,0.63,0.44}{{#1}}}
    \newcommand{\BaseNTok}[1]{\textcolor[rgb]{0.25,0.63,0.44}{{#1}}}
    \newcommand{\FloatTok}[1]{\textcolor[rgb]{0.25,0.63,0.44}{{#1}}}
    \newcommand{\CharTok}[1]{\textcolor[rgb]{0.25,0.44,0.63}{{#1}}}
    \newcommand{\StringTok}[1]{\textcolor[rgb]{0.25,0.44,0.63}{{#1}}}
    \newcommand{\CommentTok}[1]{\textcolor[rgb]{0.38,0.63,0.69}{\textit{{#1}}}}
    \newcommand{\OtherTok}[1]{\textcolor[rgb]{0.00,0.44,0.13}{{#1}}}
    \newcommand{\AlertTok}[1]{\textcolor[rgb]{1.00,0.00,0.00}{\textbf{{#1}}}}
    \newcommand{\FunctionTok}[1]{\textcolor[rgb]{0.02,0.16,0.49}{{#1}}}
    \newcommand{\RegionMarkerTok}[1]{{#1}}
    \newcommand{\ErrorTok}[1]{\textcolor[rgb]{1.00,0.00,0.00}{\textbf{{#1}}}}
    \newcommand{\NormalTok}[1]{{#1}}
    
    % Additional commands for more recent versions of Pandoc
    \newcommand{\ConstantTok}[1]{\textcolor[rgb]{0.53,0.00,0.00}{{#1}}}
    \newcommand{\SpecialCharTok}[1]{\textcolor[rgb]{0.25,0.44,0.63}{{#1}}}
    \newcommand{\VerbatimStringTok}[1]{\textcolor[rgb]{0.25,0.44,0.63}{{#1}}}
    \newcommand{\SpecialStringTok}[1]{\textcolor[rgb]{0.73,0.40,0.53}{{#1}}}
    \newcommand{\ImportTok}[1]{{#1}}
    \newcommand{\DocumentationTok}[1]{\textcolor[rgb]{0.73,0.13,0.13}{\textit{{#1}}}}
    \newcommand{\AnnotationTok}[1]{\textcolor[rgb]{0.38,0.63,0.69}{\textbf{\textit{{#1}}}}}
    \newcommand{\CommentVarTok}[1]{\textcolor[rgb]{0.38,0.63,0.69}{\textbf{\textit{{#1}}}}}
    \newcommand{\VariableTok}[1]{\textcolor[rgb]{0.10,0.09,0.49}{{#1}}}
    \newcommand{\ControlFlowTok}[1]{\textcolor[rgb]{0.00,0.44,0.13}{\textbf{{#1}}}}
    \newcommand{\OperatorTok}[1]{\textcolor[rgb]{0.40,0.40,0.40}{{#1}}}
    \newcommand{\BuiltInTok}[1]{{#1}}
    \newcommand{\ExtensionTok}[1]{{#1}}
    \newcommand{\PreprocessorTok}[1]{\textcolor[rgb]{0.74,0.48,0.00}{{#1}}}
    \newcommand{\AttributeTok}[1]{\textcolor[rgb]{0.49,0.56,0.16}{{#1}}}
    \newcommand{\InformationTok}[1]{\textcolor[rgb]{0.38,0.63,0.69}{\textbf{\textit{{#1}}}}}
    \newcommand{\WarningTok}[1]{\textcolor[rgb]{0.38,0.63,0.69}{\textbf{\textit{{#1}}}}}
    
    
    % Define a nice break command that doesn't care if a line doesn't already
    % exist.
    \def\br{\hspace*{\fill} \\* }
    % Math Jax compatibility definitions
    \def\gt{>}
    \def\lt{<}
    \let\Oldtex\TeX
    \let\Oldlatex\LaTeX
    \renewcommand{\TeX}{\textrm{\Oldtex}}
    \renewcommand{\LaTeX}{\textrm{\Oldlatex}}
    % Document parameters
    % Document title
    \title{08.is\_everything\_alright\_with\_Gauss\_theorem}
    
    
    
    
    
% Pygments definitions
\makeatletter
\def\PY@reset{\let\PY@it=\relax \let\PY@bf=\relax%
    \let\PY@ul=\relax \let\PY@tc=\relax%
    \let\PY@bc=\relax \let\PY@ff=\relax}
\def\PY@tok#1{\csname PY@tok@#1\endcsname}
\def\PY@toks#1+{\ifx\relax#1\empty\else%
    \PY@tok{#1}\expandafter\PY@toks\fi}
\def\PY@do#1{\PY@bc{\PY@tc{\PY@ul{%
    \PY@it{\PY@bf{\PY@ff{#1}}}}}}}
\def\PY#1#2{\PY@reset\PY@toks#1+\relax+\PY@do{#2}}

\@namedef{PY@tok@w}{\def\PY@tc##1{\textcolor[rgb]{0.73,0.73,0.73}{##1}}}
\@namedef{PY@tok@c}{\let\PY@it=\textit\def\PY@tc##1{\textcolor[rgb]{0.25,0.50,0.50}{##1}}}
\@namedef{PY@tok@cp}{\def\PY@tc##1{\textcolor[rgb]{0.74,0.48,0.00}{##1}}}
\@namedef{PY@tok@k}{\let\PY@bf=\textbf\def\PY@tc##1{\textcolor[rgb]{0.00,0.50,0.00}{##1}}}
\@namedef{PY@tok@kp}{\def\PY@tc##1{\textcolor[rgb]{0.00,0.50,0.00}{##1}}}
\@namedef{PY@tok@kt}{\def\PY@tc##1{\textcolor[rgb]{0.69,0.00,0.25}{##1}}}
\@namedef{PY@tok@o}{\def\PY@tc##1{\textcolor[rgb]{0.40,0.40,0.40}{##1}}}
\@namedef{PY@tok@ow}{\let\PY@bf=\textbf\def\PY@tc##1{\textcolor[rgb]{0.67,0.13,1.00}{##1}}}
\@namedef{PY@tok@nb}{\def\PY@tc##1{\textcolor[rgb]{0.00,0.50,0.00}{##1}}}
\@namedef{PY@tok@nf}{\def\PY@tc##1{\textcolor[rgb]{0.00,0.00,1.00}{##1}}}
\@namedef{PY@tok@nc}{\let\PY@bf=\textbf\def\PY@tc##1{\textcolor[rgb]{0.00,0.00,1.00}{##1}}}
\@namedef{PY@tok@nn}{\let\PY@bf=\textbf\def\PY@tc##1{\textcolor[rgb]{0.00,0.00,1.00}{##1}}}
\@namedef{PY@tok@ne}{\let\PY@bf=\textbf\def\PY@tc##1{\textcolor[rgb]{0.82,0.25,0.23}{##1}}}
\@namedef{PY@tok@nv}{\def\PY@tc##1{\textcolor[rgb]{0.10,0.09,0.49}{##1}}}
\@namedef{PY@tok@no}{\def\PY@tc##1{\textcolor[rgb]{0.53,0.00,0.00}{##1}}}
\@namedef{PY@tok@nl}{\def\PY@tc##1{\textcolor[rgb]{0.63,0.63,0.00}{##1}}}
\@namedef{PY@tok@ni}{\let\PY@bf=\textbf\def\PY@tc##1{\textcolor[rgb]{0.60,0.60,0.60}{##1}}}
\@namedef{PY@tok@na}{\def\PY@tc##1{\textcolor[rgb]{0.49,0.56,0.16}{##1}}}
\@namedef{PY@tok@nt}{\let\PY@bf=\textbf\def\PY@tc##1{\textcolor[rgb]{0.00,0.50,0.00}{##1}}}
\@namedef{PY@tok@nd}{\def\PY@tc##1{\textcolor[rgb]{0.67,0.13,1.00}{##1}}}
\@namedef{PY@tok@s}{\def\PY@tc##1{\textcolor[rgb]{0.73,0.13,0.13}{##1}}}
\@namedef{PY@tok@sd}{\let\PY@it=\textit\def\PY@tc##1{\textcolor[rgb]{0.73,0.13,0.13}{##1}}}
\@namedef{PY@tok@si}{\let\PY@bf=\textbf\def\PY@tc##1{\textcolor[rgb]{0.73,0.40,0.53}{##1}}}
\@namedef{PY@tok@se}{\let\PY@bf=\textbf\def\PY@tc##1{\textcolor[rgb]{0.73,0.40,0.13}{##1}}}
\@namedef{PY@tok@sr}{\def\PY@tc##1{\textcolor[rgb]{0.73,0.40,0.53}{##1}}}
\@namedef{PY@tok@ss}{\def\PY@tc##1{\textcolor[rgb]{0.10,0.09,0.49}{##1}}}
\@namedef{PY@tok@sx}{\def\PY@tc##1{\textcolor[rgb]{0.00,0.50,0.00}{##1}}}
\@namedef{PY@tok@m}{\def\PY@tc##1{\textcolor[rgb]{0.40,0.40,0.40}{##1}}}
\@namedef{PY@tok@gh}{\let\PY@bf=\textbf\def\PY@tc##1{\textcolor[rgb]{0.00,0.00,0.50}{##1}}}
\@namedef{PY@tok@gu}{\let\PY@bf=\textbf\def\PY@tc##1{\textcolor[rgb]{0.50,0.00,0.50}{##1}}}
\@namedef{PY@tok@gd}{\def\PY@tc##1{\textcolor[rgb]{0.63,0.00,0.00}{##1}}}
\@namedef{PY@tok@gi}{\def\PY@tc##1{\textcolor[rgb]{0.00,0.63,0.00}{##1}}}
\@namedef{PY@tok@gr}{\def\PY@tc##1{\textcolor[rgb]{1.00,0.00,0.00}{##1}}}
\@namedef{PY@tok@ge}{\let\PY@it=\textit}
\@namedef{PY@tok@gs}{\let\PY@bf=\textbf}
\@namedef{PY@tok@gp}{\let\PY@bf=\textbf\def\PY@tc##1{\textcolor[rgb]{0.00,0.00,0.50}{##1}}}
\@namedef{PY@tok@go}{\def\PY@tc##1{\textcolor[rgb]{0.53,0.53,0.53}{##1}}}
\@namedef{PY@tok@gt}{\def\PY@tc##1{\textcolor[rgb]{0.00,0.27,0.87}{##1}}}
\@namedef{PY@tok@err}{\def\PY@bc##1{{\setlength{\fboxsep}{\string -\fboxrule}\fcolorbox[rgb]{1.00,0.00,0.00}{1,1,1}{\strut ##1}}}}
\@namedef{PY@tok@kc}{\let\PY@bf=\textbf\def\PY@tc##1{\textcolor[rgb]{0.00,0.50,0.00}{##1}}}
\@namedef{PY@tok@kd}{\let\PY@bf=\textbf\def\PY@tc##1{\textcolor[rgb]{0.00,0.50,0.00}{##1}}}
\@namedef{PY@tok@kn}{\let\PY@bf=\textbf\def\PY@tc##1{\textcolor[rgb]{0.00,0.50,0.00}{##1}}}
\@namedef{PY@tok@kr}{\let\PY@bf=\textbf\def\PY@tc##1{\textcolor[rgb]{0.00,0.50,0.00}{##1}}}
\@namedef{PY@tok@bp}{\def\PY@tc##1{\textcolor[rgb]{0.00,0.50,0.00}{##1}}}
\@namedef{PY@tok@fm}{\def\PY@tc##1{\textcolor[rgb]{0.00,0.00,1.00}{##1}}}
\@namedef{PY@tok@vc}{\def\PY@tc##1{\textcolor[rgb]{0.10,0.09,0.49}{##1}}}
\@namedef{PY@tok@vg}{\def\PY@tc##1{\textcolor[rgb]{0.10,0.09,0.49}{##1}}}
\@namedef{PY@tok@vi}{\def\PY@tc##1{\textcolor[rgb]{0.10,0.09,0.49}{##1}}}
\@namedef{PY@tok@vm}{\def\PY@tc##1{\textcolor[rgb]{0.10,0.09,0.49}{##1}}}
\@namedef{PY@tok@sa}{\def\PY@tc##1{\textcolor[rgb]{0.73,0.13,0.13}{##1}}}
\@namedef{PY@tok@sb}{\def\PY@tc##1{\textcolor[rgb]{0.73,0.13,0.13}{##1}}}
\@namedef{PY@tok@sc}{\def\PY@tc##1{\textcolor[rgb]{0.73,0.13,0.13}{##1}}}
\@namedef{PY@tok@dl}{\def\PY@tc##1{\textcolor[rgb]{0.73,0.13,0.13}{##1}}}
\@namedef{PY@tok@s2}{\def\PY@tc##1{\textcolor[rgb]{0.73,0.13,0.13}{##1}}}
\@namedef{PY@tok@sh}{\def\PY@tc##1{\textcolor[rgb]{0.73,0.13,0.13}{##1}}}
\@namedef{PY@tok@s1}{\def\PY@tc##1{\textcolor[rgb]{0.73,0.13,0.13}{##1}}}
\@namedef{PY@tok@mb}{\def\PY@tc##1{\textcolor[rgb]{0.40,0.40,0.40}{##1}}}
\@namedef{PY@tok@mf}{\def\PY@tc##1{\textcolor[rgb]{0.40,0.40,0.40}{##1}}}
\@namedef{PY@tok@mh}{\def\PY@tc##1{\textcolor[rgb]{0.40,0.40,0.40}{##1}}}
\@namedef{PY@tok@mi}{\def\PY@tc##1{\textcolor[rgb]{0.40,0.40,0.40}{##1}}}
\@namedef{PY@tok@il}{\def\PY@tc##1{\textcolor[rgb]{0.40,0.40,0.40}{##1}}}
\@namedef{PY@tok@mo}{\def\PY@tc##1{\textcolor[rgb]{0.40,0.40,0.40}{##1}}}
\@namedef{PY@tok@ch}{\let\PY@it=\textit\def\PY@tc##1{\textcolor[rgb]{0.25,0.50,0.50}{##1}}}
\@namedef{PY@tok@cm}{\let\PY@it=\textit\def\PY@tc##1{\textcolor[rgb]{0.25,0.50,0.50}{##1}}}
\@namedef{PY@tok@cpf}{\let\PY@it=\textit\def\PY@tc##1{\textcolor[rgb]{0.25,0.50,0.50}{##1}}}
\@namedef{PY@tok@c1}{\let\PY@it=\textit\def\PY@tc##1{\textcolor[rgb]{0.25,0.50,0.50}{##1}}}
\@namedef{PY@tok@cs}{\let\PY@it=\textit\def\PY@tc##1{\textcolor[rgb]{0.25,0.50,0.50}{##1}}}

\def\PYZbs{\char`\\}
\def\PYZus{\char`\_}
\def\PYZob{\char`\{}
\def\PYZcb{\char`\}}
\def\PYZca{\char`\^}
\def\PYZam{\char`\&}
\def\PYZlt{\char`\<}
\def\PYZgt{\char`\>}
\def\PYZsh{\char`\#}
\def\PYZpc{\char`\%}
\def\PYZdl{\char`\$}
\def\PYZhy{\char`\-}
\def\PYZsq{\char`\'}
\def\PYZdq{\char`\"}
\def\PYZti{\char`\~}
% for compatibility with earlier versions
\def\PYZat{@}
\def\PYZlb{[}
\def\PYZrb{]}
\makeatother


    % For linebreaks inside Verbatim environment from package fancyvrb. 
    \makeatletter
        \newbox\Wrappedcontinuationbox 
        \newbox\Wrappedvisiblespacebox 
        \newcommand*\Wrappedvisiblespace {\textcolor{red}{\textvisiblespace}} 
        \newcommand*\Wrappedcontinuationsymbol {\textcolor{red}{\llap{\tiny$\m@th\hookrightarrow$}}} 
        \newcommand*\Wrappedcontinuationindent {3ex } 
        \newcommand*\Wrappedafterbreak {\kern\Wrappedcontinuationindent\copy\Wrappedcontinuationbox} 
        % Take advantage of the already applied Pygments mark-up to insert 
        % potential linebreaks for TeX processing. 
        %        {, <, #, %, $, ' and ": go to next line. 
        %        _, }, ^, &, >, - and ~: stay at end of broken line. 
        % Use of \textquotesingle for straight quote. 
        \newcommand*\Wrappedbreaksatspecials {% 
            \def\PYGZus{\discretionary{\char`\_}{\Wrappedafterbreak}{\char`\_}}% 
            \def\PYGZob{\discretionary{}{\Wrappedafterbreak\char`\{}{\char`\{}}% 
            \def\PYGZcb{\discretionary{\char`\}}{\Wrappedafterbreak}{\char`\}}}% 
            \def\PYGZca{\discretionary{\char`\^}{\Wrappedafterbreak}{\char`\^}}% 
            \def\PYGZam{\discretionary{\char`\&}{\Wrappedafterbreak}{\char`\&}}% 
            \def\PYGZlt{\discretionary{}{\Wrappedafterbreak\char`\<}{\char`\<}}% 
            \def\PYGZgt{\discretionary{\char`\>}{\Wrappedafterbreak}{\char`\>}}% 
            \def\PYGZsh{\discretionary{}{\Wrappedafterbreak\char`\#}{\char`\#}}% 
            \def\PYGZpc{\discretionary{}{\Wrappedafterbreak\char`\%}{\char`\%}}% 
            \def\PYGZdl{\discretionary{}{\Wrappedafterbreak\char`\$}{\char`\$}}% 
            \def\PYGZhy{\discretionary{\char`\-}{\Wrappedafterbreak}{\char`\-}}% 
            \def\PYGZsq{\discretionary{}{\Wrappedafterbreak\textquotesingle}{\textquotesingle}}% 
            \def\PYGZdq{\discretionary{}{\Wrappedafterbreak\char`\"}{\char`\"}}% 
            \def\PYGZti{\discretionary{\char`\~}{\Wrappedafterbreak}{\char`\~}}% 
        } 
        % Some characters . , ; ? ! / are not pygmentized. 
        % This macro makes them "active" and they will insert potential linebreaks 
        \newcommand*\Wrappedbreaksatpunct {% 
            \lccode`\~`\.\lowercase{\def~}{\discretionary{\hbox{\char`\.}}{\Wrappedafterbreak}{\hbox{\char`\.}}}% 
            \lccode`\~`\,\lowercase{\def~}{\discretionary{\hbox{\char`\,}}{\Wrappedafterbreak}{\hbox{\char`\,}}}% 
            \lccode`\~`\;\lowercase{\def~}{\discretionary{\hbox{\char`\;}}{\Wrappedafterbreak}{\hbox{\char`\;}}}% 
            \lccode`\~`\:\lowercase{\def~}{\discretionary{\hbox{\char`\:}}{\Wrappedafterbreak}{\hbox{\char`\:}}}% 
            \lccode`\~`\?\lowercase{\def~}{\discretionary{\hbox{\char`\?}}{\Wrappedafterbreak}{\hbox{\char`\?}}}% 
            \lccode`\~`\!\lowercase{\def~}{\discretionary{\hbox{\char`\!}}{\Wrappedafterbreak}{\hbox{\char`\!}}}% 
            \lccode`\~`\/\lowercase{\def~}{\discretionary{\hbox{\char`\/}}{\Wrappedafterbreak}{\hbox{\char`\/}}}% 
            \catcode`\.\active
            \catcode`\,\active 
            \catcode`\;\active
            \catcode`\:\active
            \catcode`\?\active
            \catcode`\!\active
            \catcode`\/\active 
            \lccode`\~`\~ 	
        }
    \makeatother

    \let\OriginalVerbatim=\Verbatim
    \makeatletter
    \renewcommand{\Verbatim}[1][1]{%
        %\parskip\z@skip
        \sbox\Wrappedcontinuationbox {\Wrappedcontinuationsymbol}%
        \sbox\Wrappedvisiblespacebox {\FV@SetupFont\Wrappedvisiblespace}%
        \def\FancyVerbFormatLine ##1{\hsize\linewidth
            \vtop{\raggedright\hyphenpenalty\z@\exhyphenpenalty\z@
                \doublehyphendemerits\z@\finalhyphendemerits\z@
                \strut ##1\strut}%
        }%
        % If the linebreak is at a space, the latter will be displayed as visible
        % space at end of first line, and a continuation symbol starts next line.
        % Stretch/shrink are however usually zero for typewriter font.
        \def\FV@Space {%
            \nobreak\hskip\z@ plus\fontdimen3\font minus\fontdimen4\font
            \discretionary{\copy\Wrappedvisiblespacebox}{\Wrappedafterbreak}
            {\kern\fontdimen2\font}%
        }%
        
        % Allow breaks at special characters using \PYG... macros.
        \Wrappedbreaksatspecials
        % Breaks at punctuation characters . , ; ? ! and / need catcode=\active 	
        \OriginalVerbatim[#1,codes*=\Wrappedbreaksatpunct]%
    }
    \makeatother

    % Exact colors from NB
    \definecolor{incolor}{HTML}{303F9F}
    \definecolor{outcolor}{HTML}{D84315}
    \definecolor{cellborder}{HTML}{CFCFCF}
    \definecolor{cellbackground}{HTML}{F7F7F7}
    
    % prompt
    \makeatletter
    \newcommand{\boxspacing}{\kern\kvtcb@left@rule\kern\kvtcb@boxsep}
    \makeatother
    \newcommand{\prompt}[4]{
        {\ttfamily\llap{{\color{#2}[#3]:\hspace{3pt}#4}}\vspace{-\baselineskip}}
    }
    

    
    % Prevent overflowing lines due to hard-to-break entities
    \sloppy 
    % Setup hyperref package
    \hypersetup{
      breaklinks=true,  % so long urls are correctly broken across lines
      colorlinks=true,
      urlcolor=urlcolor,
      linkcolor=linkcolor,
      citecolor=citecolor,
      }
    % Slightly bigger margins than the latex defaults
    
    \geometry{verbose,tmargin=1in,bmargin=1in,lmargin=1in,rmargin=1in}
    
    

\begin{document}
    
    \maketitle
    
    

    
    \section{Все ли в порядке с теоремой
Гаусса?}\label{ux432ux441ux435-ux43bux438-ux432-ux43fux43eux440ux44fux434ux43aux435-ux441-ux442ux435ux43eux440ux435ux43cux43eux439-ux433ux430ux443ux441ux441ux430}

Сумма двух компонент электрического поля
\[\overrightarrow{E}=-\nabla \varphi -\frac{1}{c}\frac{\partial \overrightarrow{A}}{\partial t}=\frac{e}{{{K}^{3}}}\left[ \left( \overrightarrow{R}-R\frac{\overrightarrow{v}}{c} \right)\left( 1+\frac{\overrightarrow{a}\cdot \overrightarrow{R}}{{{c}^{2}}}-\frac{\overrightarrow{v}\cdot \overrightarrow{v}}{{{c}^{2}}} \right)-KR\frac{\overrightarrow{a}}{{{c}^{2}}} \right]\]
\[\overrightarrow{E}=\frac{e}{{{K}^{3}}}\left( \overrightarrow{R}-R\frac{\overrightarrow{v}}{c} \right)\left( 1+\frac{\overrightarrow{a}\cdot \overrightarrow{R}}{{{c}^{2}}}-\frac{\overrightarrow{v}\cdot \overrightarrow{v}}{{{c}^{2}}} \right)-\frac{e}{{{K}^{2}}}R\frac{\overrightarrow{a}}{{{c}^{2}}}\]
Пусть заряд в запаздывающий момент времени \({{t}_{0}}'\) находится в
начале координат. Вообразим сферу радиуса \({{R}_{0}}\) с центром в
начале координат и вычислим нормальную компоненту электрического поля в
момент \[{{t}_{0}}={{t}_{0}}'+\frac{{{R}_{0}}}{c}\] на поверхности этой
сферы. Для этого вектор \(\overrightarrow{E}\) необходимо скалярно
умножить на вектор нормали
\(\overrightarrow{n}=\frac{\overrightarrow{{{R}_{0}}}}{{{R}_{0}}}\) ,
где \(\overrightarrow{{{R}_{0}}}\) есть вектор проведённый из начала
координат в точку поверхности сферы. При данных условиях, поскольку для
всей поверхности сферы запаздывающий момент одинаков и заряд в этот
запаздывающий момент находится в начале координат, то \(R={{R}_{0}}\) и
\(\overrightarrow{R}=\overrightarrow{{{R}_{0}}}\). Отсюда
\({{E}_{n}}=\overrightarrow{E}\frac{\overrightarrow{R}}{R}=\frac{e}{{{K}^{3}}}\left( \overrightarrow{R}\frac{\overrightarrow{R}}{R}-R\frac{\overrightarrow{v}}{c}\frac{\overrightarrow{R}}{R} \right)\left( 1+\frac{\overrightarrow{a}\cdot \overrightarrow{R}}{{{c}^{2}}}-\frac{\overrightarrow{v}\cdot \overrightarrow{v}}{{{c}^{2}}} \right)-\frac{e}{{{K}^{2}}}R\frac{\overrightarrow{a}}{{{c}^{2}}}\frac{\overrightarrow{R}}{R}\)
\({{E}_{n}}=\frac{e}{{{K}^{3}}}\left( R-\frac{\overrightarrow{v}\cdot \overrightarrow{R}}{c} \right)\left( 1+\frac{\overrightarrow{a}\cdot \overrightarrow{R}}{{{c}^{2}}}-\frac{\overrightarrow{v}\cdot \overrightarrow{v}}{{{c}^{2}}} \right)-\frac{e}{{{K}^{2}}}\frac{\overrightarrow{a}\cdot \overrightarrow{R}}{{{c}^{2}}}\)
\({{E}_{n}}=\frac{e}{{{K}^{2}}}\left( 1+\frac{\overrightarrow{a}\cdot \overrightarrow{R}}{{{c}^{2}}}-\frac{\overrightarrow{v}\cdot \overrightarrow{v}}{{{c}^{2}}} \right)-\frac{e}{{{K}^{2}}}\frac{\overrightarrow{a}\cdot \overrightarrow{R}}{{{c}^{2}}}\)
\({{E}_{n}}=\frac{e}{{{K}^{2}}}\left( 1-\frac{\overrightarrow{v}\cdot \overrightarrow{v}}{{{c}^{2}}} \right)=\frac{e}{{{\left( R-\frac{\overrightarrow{v}\cdot \overrightarrow{R}}{c} \right)}^{2}}}\left( 1-\frac{\overrightarrow{v}\cdot \overrightarrow{v}}{{{c}^{2}}} \right)\)
Интересно, что ускорение заряда не входит в формулу для нормальной
компоненты электрического поля. Нормальная компонента электрического
поля зависит только от скорости заряда в момент \({{t}_{0}}'\). Выберем
направление скорости заряда в этот момент в качестве оси и введём
сферическую систему координат. В этой системе
\(\overrightarrow{v}\cdot \overrightarrow{R}=vR\cos \left( \theta \right)\)
\({{E}_{n}}=e\frac{\left( 1-\frac{{{v}^{2}}}{{{c}^{2}}} \right)}{{{\left( R-\frac{vR\cos \left( \theta \right)}{c} \right)}^{2}}}=\frac{e}{{{R}^{2}}}\frac{\left( 1-\frac{{{v}^{2}}}{{{c}^{2}}} \right)}{{{\left( 1-\frac{v}{c}\cos \left( \theta \right) \right)}^{2}}}\)
Интеграл по поверхности сферы
\[\iint{{{E}_{n}}{{R}^{2}}}\sin \left( \theta  \right)d\theta d\varphi =e\int\limits_{\theta =0}^{\pi }{\int\limits_{\varphi =0}^{2\pi }{\frac{\left( 1-\frac{{{v}^{2}}}{{{c}^{2}}} \right)}{{{\left( 1-\frac{v}{c}\cos \left( \theta  \right) \right)}^{2}}}}}\sin \left( \theta  \right)d\theta d\varphi =2\pi e\int\limits_{\theta =0}^{\pi }{\frac{\left( 1-\frac{{{v}^{2}}}{{{c}^{2}}} \right)}{{{\left( 1-\frac{v}{c}\cos \left( \theta  \right) \right)}^{2}}}}\sin \left( \theta  \right)d\theta =4\pi e\]
В данном рассмотренном выше случае соответствует теореме Гаусса. Но
возникает вопрос, будет соответствовать ли теореме Гаусса этот интеграл
по той же поверхности, но в другие моменты времени? Для простоты
рассмотрим практически важный случай, когда векторы ускорения и скорости
заряда сонаправлены. Уравнение движения заряда для этого случая
\[{{z}_{q}}={{z}_{0}}+{{v}_{0}}\left( t'-{{t}_{0}}' \right)+\frac{{{a}_{0}}{{\left( t'-{{t}_{0}}' \right)}^{2}}}{2}\]
\[{{v}_{q}}={{v}_{0}}+{{a}_{0}}\left( t'-{{t}_{0}}' \right)\] Расчёт
запаздывающего момента \(t'\) производится с помощью разрешения
уравнения вида (это уравнение записано с помощью теоремы косинусов)
\[{{R}^{2}}={{c}^{2}}{{\left( t-t' \right)}^{2}}={{R}_{0}}^{2}-2{{R}_{0}}{{z}_{q}}\cos \theta +{{z}_{q}}^{2}\]
Нормальная компонента электрического поля
\({{E}_{n}}=\overrightarrow{E}\frac{\overrightarrow{{{R}_{0}}}}{{{R}_{0}}}=\frac{e}{{{K}^{3}}}\left( \overrightarrow{R}\frac{\overrightarrow{{{R}_{0}}}}{{{R}_{0}}}-R\frac{\overrightarrow{v}}{c}\frac{\overrightarrow{{{R}_{0}}}}{{{R}_{0}}} \right)\left( 1+\frac{\overrightarrow{a}\cdot \overrightarrow{R}}{{{c}^{2}}}-\frac{\overrightarrow{v}\cdot \overrightarrow{v}}{{{c}^{2}}} \right)-\frac{e}{{{K}^{2}}}R\frac{\overrightarrow{a}}{{{c}^{2}}}\frac{\overrightarrow{{{R}_{0}}}}{{{R}_{0}}}\)
Для упрощения аналитического вывода формулы запаздывающего момента можно
применить два приближения: приближение малых скоростей \(v\to 0\) и
приближение малых ускорений \(a\to 0\) В приближении малых скоростей при
\({{z}_{0}}=0\) уравнение для вычисления запаздывающего момента
принимает вид
\[{{c}^{2}}{{\left( t-t' \right)}^{2}}={{R}_{0}}^{2}-2{{R}_{0}}\frac{{{a}_{0}}{{\left( t'-{{t}_{0}}' \right)}^{2}}}{2}\cos \theta +{{\left( \frac{{{a}_{0}}{{\left( t'-{{t}_{0}}' \right)}^{2}}}{2} \right)}^{2}}\]
А в приближении малых ускорений при \({{z}_{0}}=0\)
\[{{c}^{2}}{{\left( t-t' \right)}^{2}}={{R}_{0}}^{2}-2{{R}_{0}}{{v}_{0}}\left( t'-{{t}_{0}}' \right)\cos \theta +{{\left( {{v}_{0}}\left( t'-{{t}_{0}}' \right) \right)}^{2}}\]
В приближении малых ускорений аналитически решать все же проще
\[t'=\frac{-{{R}_{0}}{{v}_{0}}cos\left( \theta  \right)+{{c}^{2}}t-{{t}_{0}}'{{v}_{0}}^{2}-\sqrt{cos{{\left( \theta  \right)}^{2}}{{R}_{0}}^{2}{{v}_{0}}^{2}-2cos\left( \theta  \right){{R}_{0}}{{c}^{2}}t{{v}_{0}}+2cos\left( \theta  \right){{R}_{0}}{{c}^{2}}{{t}_{0}}'{{v}_{0}}+{{c}^{2}}{{t}^{2}}{{v}_{0}}^{2}-2{{c}^{2}}t\cdot {{t}_{0}}'{{v}_{0}}^{2}+{{c}^{2}}{{\left( {{t}_{0}}' \right)}^{2}}{{v}_{0}}^{2}+{{R}_{0}}^{2}{{c}^{2}}-{{R}_{0}}^{2}{{v}_{0}}^{2}}}{({{c}^{2}}-{{v}_{0}}^{2})}\]
\[t'=\frac{-{{R}_{0}}{{v}_{0}}cos\left( \theta  \right)+{{c}^{2}}t-{{t}_{0}}'{{v}_{0}}^{2}-\sqrt{cos{{\left( \theta  \right)}^{2}}{{R}_{0}}^{2}{{v}_{0}}^{2}-2cos\left( \theta  \right){{R}_{0}}{{v}_{0}}\left( t-{{t}_{0}}' \right){{c}^{2}}+{{\left( ct{{v}_{0}} \right)}^{2}}-2{{c}^{2}}t\left( {{t}_{0}}' \right){{v}_{0}}^{2}+{{\left( c\left( {{t}_{0}}' \right){{v}_{0}} \right)}^{2}}+{{R}_{0}}^{2}\left( {{c}^{2}}-{{v}_{0}}^{2} \right)}}{({{c}^{2}}-{{v}_{0}}^{2})}\]
\[t'=\frac{-{{R}_{0}}{{v}_{0}}cos\left( \theta  \right)+{{c}^{2}}t-{{t}_{0}}'{{v}_{0}}^{2}-\sqrt{{{\left( cos\left( \theta  \right){{R}_{0}}{{v}_{0}} \right)}^{2}}-2cos\left( \theta  \right){{R}_{0}}{{v}_{0}}\left( t-{{t}_{0}}' \right){{c}^{2}}+{{c}^{2}}{{v}_{0}}^{2}{{\left( t-{{t}_{0}}' \right)}^{2}}+{{R}_{0}}^{2}\left( {{c}^{2}}-{{v}_{0}}^{2} \right)}}{({{c}^{2}}-{{v}_{0}}^{2})}\]
Откуда
\[R=c\left( t-t' \right)=c\left( t-\frac{-{{R}_{0}}{{v}_{0}}cos\left( \theta  \right)+{{c}^{2}}t-{{t}_{0}}'{{v}_{0}}^{2}-\sqrt{{{\left( cos\left( \theta  \right){{R}_{0}}{{v}_{0}} \right)}^{2}}-2cos\left( \theta  \right){{R}_{0}}{{v}_{0}}\left( t-{{t}_{0}}' \right){{c}^{2}}+{{c}^{2}}{{v}_{0}}^{2}{{\left( t-{{t}_{0}}' \right)}^{2}}+{{R}_{0}}^{2}\left( {{c}^{2}}-{{v}_{0}}^{2} \right)}}{({{c}^{2}}-{{v}_{0}}^{2})} \right)\]
Нормальная компонента электрического поля в приближении малых ускорений
\(a\to 0\)
\({{E}_{n}}=\frac{e}{{{K}^{3}}}\left( \overrightarrow{R}\frac{\overrightarrow{{{R}_{0}}}}{{{R}_{0}}}-R\frac{\overrightarrow{v}}{c}\frac{\overrightarrow{{{R}_{0}}}}{{{R}_{0}}} \right)\left( 1-\frac{\overrightarrow{v}\cdot \overrightarrow{v}}{{{c}^{2}}} \right)\)
Теперь необходимо определить радиус векторы. В декартовой системе
\[\overrightarrow{R}=\overrightarrow{i}\left( x \right)+\overrightarrow{j}\left( y \right)+\overrightarrow{k}\left( z-{{z}_{q}} \right)\]
\[\overrightarrow{{{R}_{0}}}=\overrightarrow{i}\left( x \right)+\overrightarrow{j}\left( y \right)+\overrightarrow{k}\left( z \right)\]
где \(x={{R}_{0}}\sin \left( \theta \right)\cos \left( \varphi \right)\)
\(y={{R}_{0}}\sin \left( \theta \right)\sin \left( \varphi \right)\)
\(z={{R}_{0}}\cos \left( \theta \right)\) а
\[{{z}_{q}}={{z}_{0}}+{{v}_{0}}\left( t'-{{t}_{0}}' \right)\]
Следовательно, с использованием сферических координат эти векторы
запишутся
\[\overrightarrow{R}=\overrightarrow{i}{{R}_{0}}\sin \left( \theta  \right)\cos \left( \varphi  \right)+\overrightarrow{j}{{R}_{0}}\sin \left( \theta  \right)\sin \left( \varphi  \right)+\overrightarrow{k}\left( {{R}_{0}}\cos \left( \theta  \right)-{{z}_{q}} \right)\]
\[\overrightarrow{{{R}_{0}}}=\overrightarrow{i}{{R}_{0}}\sin \left( \theta  \right)\cos \left( \varphi  \right)+\overrightarrow{j}{{R}_{0}}\sin \left( \theta  \right)\sin \left( \varphi  \right)+\overrightarrow{k}{{R}_{0}}\cos \left( \theta  \right)\]
Скалярные произведения
\[\overrightarrow{R}\frac{\overrightarrow{{{R}_{0}}}}{{{R}_{0}}}={{R}_{0}}\sin {{\left( \theta  \right)}^{2}}\cos {{\left( \varphi  \right)}^{2}}+{{R}_{0}}\sin {{\left( \theta  \right)}^{2}}\sin {{\left( \varphi  \right)}^{2}}+{{R}_{0}}\cos {{\left( \theta  \right)}^{2}}-{{z}_{q}}\cos \left( \theta  \right)\]
\[\overrightarrow{R}\frac{\overrightarrow{{{R}_{0}}}}{{{R}_{0}}}={{R}_{0}}\sin {{\left( \theta  \right)}^{2}}+{{R}_{0}}\cos {{\left( \theta  \right)}^{2}}-{{z}_{q}}\cos \left( \theta  \right)\]
\[\overrightarrow{R}\frac{\overrightarrow{{{R}_{0}}}}{{{R}_{0}}}={{R}_{0}}-{{z}_{q}}\cos \left( \theta  \right)\]
Второе скалярное произведение, учитывая что
\[\overrightarrow{v}={{v}_{q}}\overrightarrow{k}\]
\(R\frac{\overrightarrow{v}}{c}\frac{\overrightarrow{{{R}_{0}}}}{{{R}_{0}}}=R\frac{{{v}_{q}}}{c}\cos \left( \theta \right)=c\left( t-t' \right)\frac{{{v}_{q}}}{c}\cos \left( \theta \right)=\left( t-t' \right){{v}_{q}}\cos \left( \theta \right)\)
Сумма скалярных произведений \[\begin{align}
  & \left( \overrightarrow{R}\frac{\overrightarrow{{{R}_{0}}}}{{{R}_{0}}}-R\frac{\overrightarrow{v}}{c}\frac{\overrightarrow{{{R}_{0}}}}{{{R}_{0}}} \right)=\left\{ {{R}_{0}}-{{z}_{q}}\cos \left( \theta  \right) \right\}-\left\{ \left( t-t' \right){{v}_{q}}\cos \left( \theta  \right) \right\}= \\ 
 & ={{R}_{0}}-\left( {{z}_{q}}\cos \left( \theta  \right)+\left( t-t' \right){{v}_{q}}\cos \left( \theta  \right) \right)= \\ 
 & ={{R}_{0}}-\left( {{z}_{q}}+\left( t-t' \right){{v}_{q}} \right)\cos \left( \theta  \right) \\ 
\end{align}\] радиус Лиенара Вихерта
\(\overrightarrow{v}\cdot \overrightarrow{R}={{v}_{q}}\left( {{R}_{0}}\cos \left( \theta \right)-{{z}_{q}} \right)\)
\[K=R-\frac{{{v}_{q}}\left( {{R}_{0}}\cos \left( \theta  \right)-{{z}_{q}} \right)}{c}=c\left( t-t' \right)-\frac{{{v}_{q}}\left( {{R}_{0}}\cos \left( \theta  \right)-{{z}_{q}} \right)}{c}\]
Нормальная компонента поля
\({{E}_{n}}=\frac{e}{{{K}^{3}}}\left( \overrightarrow{R}\frac{\overrightarrow{{{R}_{0}}}}{{{R}_{0}}}-R\frac{\overrightarrow{v}}{c}\frac{\overrightarrow{{{R}_{0}}}}{{{R}_{0}}} \right)\left( 1-\frac{\overrightarrow{v}\cdot \overrightarrow{v}}{{{c}^{2}}} \right)\)
\[{{E}_{n}}=\frac{e}{{{K}^{3}}}\left( {{R}_{0}}-{{z}_{q}}\cos \left( \theta  \right)-\left( t-t' \right){{v}_{q}}\cos \left( \theta  \right) \right)\left( 1-\frac{{{v}_{q}}^{2}}{{{c}^{2}}} \right)\]
\[{{E}_{n}}=\frac{e}{{{K}^{3}}}\left( {{R}_{0}}-\left( {{z}_{q}}+\left( t-t' \right){{v}_{q}} \right)\cos \left( \theta  \right) \right)\left( 1-\frac{{{v}_{q}}^{2}}{{{c}^{2}}} \right)\]
Интеграл по поверхности сферы
\[\iint{{{E}_{n}}{{R}^{2}}}\sin \left( \theta  \right)d\theta d\varphi =2\pi \int\limits_{\theta =0}^{\pi }{{{E}_{n}}{{R}_{0}}^{2}}\sin \left( \theta  \right)d\theta =2\pi e\int\limits_{\theta =0}^{\pi }{\frac{{{R}_{0}}-\left( {{z}_{q}}+\left( t-t' \right){{v}_{q}} \right)\cos \left( \theta  \right)}{{{K}^{3}}}\left( 1-\frac{{{v}_{0}}^{2}}{{{c}^{2}}} \right){{R}_{0}}^{2}}\sin \left( \theta  \right)\]
Численный расчёт этого интеграла при z\_\_0 := 0.5 R0 := 2 tau\_\_0 = 0,
v\_\_0 = 0.9, c = 1

Следовательно, для равномерно движущегося заряда теорема Гаусса
выполняется в случае использования потенциалов Лиенара Вихерта

Более общий случай (векторы ускорения и скорости заряда сонаправлены).
\[z\left( t',{{v}_{0}},{{a}_{0}} \right)={{z}_{0}}+{{v}_{0}}\left( t'-{{t}_{0}}' \right)+\frac{{{a}_{0}}{{\left( t'-{{t}_{0}}' \right)}^{2}}}{2}\]
\[v\left( t',{{v}_{0}},{{a}_{0}} \right)={{v}_{0}}+{{a}_{0}}\left( t'-{{t}_{0}}' \right)\]
Расчёт запаздывающего момента \(t'\) производится с помощью разрешения
уравнения вида (это уравнение записано с помощью теоремы косинусов)
\[{{R}^{2}}={{c}^{2}}{{\left( t-t' \right)}^{2}}={{R}_{0}}^{2}-2{{R}_{0}}z\left( t',{{v}_{0}},{{a}_{0}} \right)\cos \theta +z{{\left( t',{{v}_{0}},{{a}_{0}} \right)}^{2}}\]
Нормальная компонента электрического поля
\({{E}_{n}}=\overrightarrow{E}\frac{\overrightarrow{{{R}_{0}}}}{{{R}_{0}}}=\frac{e}{{{K}^{3}}}\left( \overrightarrow{R}\frac{\overrightarrow{{{R}_{0}}}}{{{R}_{0}}}-R\frac{\overrightarrow{v}}{c}\frac{\overrightarrow{{{R}_{0}}}}{{{R}_{0}}} \right)\left( 1+\frac{\overrightarrow{a}\cdot \overrightarrow{R}}{{{c}^{2}}}-\frac{\overrightarrow{v}\cdot \overrightarrow{v}}{{{c}^{2}}} \right)-\frac{e}{{{K}^{2}}}R\frac{\overrightarrow{a}}{{{c}^{2}}}\frac{\overrightarrow{{{R}_{0}}}}{{{R}_{0}}}\)
с использованием сферических координат векторы запишутся
\[\overrightarrow{R}=\overrightarrow{i}{{R}_{0}}\sin \left( \theta  \right)\cos \left( \varphi  \right)+\overrightarrow{j}{{R}_{0}}\sin \left( \theta  \right)\sin \left( \varphi  \right)+\overrightarrow{k}\left( {{R}_{0}}\cos \left( \theta  \right)-{{z}_{q}}\left( t' \right) \right)\]
\[\overrightarrow{{{R}_{0}}}=\overrightarrow{i}{{R}_{0}}\sin \left( \theta  \right)\cos \left( \varphi  \right)+\overrightarrow{j}{{R}_{0}}\sin \left( \theta  \right)\sin \left( \varphi  \right)+\overrightarrow{k}{{R}_{0}}\cos \left( \theta  \right)\]
Скалярные произведения

\[\overrightarrow{R}\frac{\overrightarrow{{{R}_{0}}}}{{{R}_{0}}}={{R}_{0}}\sin {{\left( \theta  \right)}^{2}}\cos {{\left( \varphi  \right)}^{2}}+{{R}_{0}}\sin {{\left( \theta  \right)}^{2}}\sin {{\left( \varphi  \right)}^{2}}+{{R}_{0}}\cos {{\left( \theta  \right)}^{2}}-{{z}_{q}}\left( t' \right)\cos \left( \theta  \right)\]
\[\overrightarrow{R}\frac{\overrightarrow{{{R}_{0}}}}{{{R}_{0}}}={{R}_{0}}\sin {{\left( \theta  \right)}^{2}}+{{R}_{0}}\cos {{\left( \theta  \right)}^{2}}-{{z}_{q}}\left( t' \right)\cos \left( \theta  \right)\]
\[\overrightarrow{R}\frac{\overrightarrow{{{R}_{0}}}}{{{R}_{0}}}={{R}_{0}}-{{z}_{q}}\left( t' \right)\cos \left( \theta  \right)\]
Второе скалярное произведение, учитывая что
\[\overrightarrow{v}={{v}_{q}}\left( t' \right)\overrightarrow{k}\]
\[R\frac{\overrightarrow{v}}{c}\frac{\overrightarrow{{{R}_{0}}}}{{{R}_{0}}}=R\frac{{{v}_{q}}\left( t' \right)}{c}\cos \left( \theta  \right)=c\left( t-t' \right)\frac{{{v}_{q}}\left( t' \right)}{c}\cos \left( \theta  \right)=\left( t-t' \right){{v}_{q}}\left( t' \right)\cos \left( \theta  \right)\]
Сумма скалярных произведений
\(\left( \overrightarrow{R}\frac{\overrightarrow{{{R}_{0}}}}{{{R}_{0}}}-R\frac{\overrightarrow{v}}{c}\frac{\overrightarrow{{{R}_{0}}}}{{{R}_{0}}} \right)={{R}_{0}}-{{z}_{q}}\left( t' \right)\cos \left( \theta \right)-\left( t-t' \right){{v}_{q}}\left( t' \right)\cos \left( \theta \right)\)
Ещё
\(\frac{\overrightarrow{a}\cdot \overrightarrow{R}}{{{c}^{2}}}={{a}_{q}}\left( t' \right)\frac{\left( {{R}_{0}}\cos \left( \theta \right)-{{z}_{q}}\left( t' \right) \right)}{{{c}^{2}}}\)
И
\(\frac{e}{{{K}^{2}}}R\frac{\overrightarrow{a}}{{{c}^{2}}}\frac{\overrightarrow{{{R}_{0}}}}{{{R}_{0}}}=\frac{e}{{{K}^{2}}}R\frac{{{a}_{q}}\left( t' \right)}{{{c}^{2}}}\frac{{{R}_{0}}\cos \left( \theta \right)}{{{R}_{0}}}=\frac{e}{{{K}^{2}}}c\left( t-t' \right)\frac{{{a}_{q}}\left( t' \right)}{{{c}^{2}}}\cos \left( \theta \right)=\frac{e}{{{K}^{2}}}\left( t-t' \right)\frac{{{a}_{q}}\left( t' \right)}{c}\cos \left( \theta \right)\)
Теперь нормальная компонента поля
\({{E}_{n}}=e\frac{{{R}_{0}}-{{z}_{q}}\left( t' \right)\cos \left( \theta \right)-\left( t-t' \right){{v}_{q}}\left( t' \right)\cos \left( \theta \right)}{{{K}^{3}}}\left( 1+{{a}_{q}}\left( t' \right)\frac{\left( {{R}_{0}}\cos \left( \theta \right)-{{z}_{q}}\left( t' \right) \right)}{{{c}^{2}}}-\frac{{{v}_{q}}{{\left( t' \right)}^{2}}}{{{c}^{2}}} \right)-\frac{e}{{{K}^{2}}}\left( t-t' \right)\frac{{{a}_{q}}\left( t' \right)}{c}\cos \left( \theta \right)\)
радиус Лиенара Вихерта
\(\overrightarrow{v}\cdot \overrightarrow{R}={{v}_{q}}\left( t' \right)\left( {{R}_{0}}\cos \left( \theta \right)-{{z}_{q}}\left( t' \right) \right)\)
\[\begin{align}
  & K=R-\frac{{{v}_{q}}\left( t' \right)\left( {{R}_{0}}\cos \left( \theta  \right)-{{z}_{q}}\left( t' \right) \right)}{c}= \\ 
 & =c\left( t-t' \right)-\frac{{{v}_{q}}\left( t' \right)\left( {{R}_{0}}\cos \left( \theta  \right)-{{z}_{q}}\left( t' \right) \right)}{c} \\ 
\end{align}\]

    \begin{tcolorbox}[breakable, size=fbox, boxrule=1pt, pad at break*=1mm,colback=cellbackground, colframe=cellborder]
\prompt{In}{incolor}{ }{\boxspacing}
\begin{Verbatim}[commandchars=\\\{\}]

\end{Verbatim}
\end{tcolorbox}


    % Add a bibliography block to the postdoc
    
    
    
\end{document}
